%% =================================================================
%% Chen et al. Measurement 264 (2026) 120300.
%% Reconstruction of page 4 (printed p. 4).
%% Contents: Fig. 6 (a-d: frame model size study) and intro
%% continuation (AFM limits, MBT positioning paragraph, self-proposal,
%% Section 2 start, Section 2.1 Scanning module).
%% No numbered equations on this page.
%% =================================================================

\documentclass[11pt]{article}

\usepackage[margin=0.85in]{geometry}
\usepackage{amsmath, amssymb}
\usepackage{mathrsfs}
\usepackage{xcolor}
\usepackage{fancyhdr}

\pagestyle{fancy}
\fancyhf{}
\lhead{\small Z.~Chen et al.}
\rhead{\small Measurement 264 (2026) 120300}
\cfoot{\small 4 $\to$ \arabic{page}}
\renewcommand{\headrulewidth}{0.4pt}

\setcounter{page}{1}
\setlength{\parskip}{6pt}
\setlength{\parindent}{0pt}

\begin{document}

% ---------- Fig. 6 placeholder ----------
\begin{center}
\fbox{\begin{minipage}[c][7cm][c]{0.95\textwidth}
\centering
\textbf{Fig. 6.}\ Effect of the size of the sample on the tactile
softness of materials.\\[4pt]
\textcolor{gray}{\footnotesize
\textbf{(a)} Frame model with a depth of 6 mm and different length
\& width (3D printed ABS grid).\\[3pt]
\textbf{(b)} Frame model filled with a soft surface layer of AB
glue. Labeled scanning area 8$\times$8 mm, with decreasing lengths
8.0, 6.4, 5.1, 4.0, 3.2 mm marked along each axis.\\[3pt]
\textbf{(c)} 3D imaging of compression depth distribution when the
tactile tomography system scanned the frame model with a soft
surface layer. X, Y in mm; Z 0--6 mm (color bar).\\[3pt]
\textbf{(d)} Top view of (c), showing that the central cells of
the grid have the largest compression depth and that depth decreases
as cell size shrinks.]}
\end{minipage}}
\end{center}

\vspace{0.4em}

\textbf{Fig. 6.}\ Effect of the size of the sample on the tactile
softness of materials. (a) Frame model with a depth of 6 mm and
different length \& width. (b) Frame model filled with a soft
surface layer of AB glue. (c) 3D imaging of compression depth
distribution when the tactile tomography system scanned the frame
model with a soft surface layer. (d) Top view of the 3D imaging of
compression depth distribution when the tactile tomography system
scanned the frame model with a soft surface layer.

\vspace{0.8em}

% ---------- Prose continuation ----------
analysis to the outermost surface or near-surface layers. Moreover,
they cannot handle thick samples ($> 1\,\mu$m) or soft-encapsulated
materials (such as flexible electronics with soft coatings), as
probe-sample interactions are too weak to penetrate without damaging
the surface.

Recent advancements in mechanical principle-based tomography have
begun to address the gap in internal imaging, with Mechanics Based
Tomography (MBT) emerging as a promising approach [32--34]. MBT
leverages solid mechanics equations and surface displacement
measurements to reconstruct internal material property distributions
through an optimization framework with regularization to handle
ill-posed inverse problems. Experimental validations on composite
silicone samples have demonstrated MBT's ability to visualize
subsurface inclusions by resolving stiffness contrasts, highlighting
its potential for non-destructive testing of soft materials [34].
Unlike X-CT or ultrasonic methods, mechanical principle-based
tomography avoids ionizing radiation or high-energy waves, making
it suitable for delicate samples. In addition, mechanical
principle-based tomography can map mechanical properties (such as
stiffness or elastic modulus), providing functional information
beyond structural imaging. However, existing mechanical tomography
methods face challenges with low reconstruction accuracy and
reliability, as displacement measurement noise and position
deviations of the force can both lead to a significant deterioration
in the quality of internal structure reconstruction.

In this study, we propose a novel tactile tomography system based
on mechanical principles, featuring internal 3D imaging
capabilities. The tactile tomography system is constructed from a
scanning module, a probe module, and a closed-loop feedback control
system, enabling it to sense force and control the movement of the
probe tip. This tactile tomography system can further recognize the
softness of soft matter. To simulate human tactile perception of
softness, a novel parameter of ``tactile softness of materials'' is
proposed. The tactile tomography system can achieve the function of
internal 3D imaging with a maximum detection depth of 6 mm by the
combined effects of the sample's size, thickness, and the applied
force. The performance of the tactile tomography system is
characterized by its scanning resolution, scanning accuracy, and the
trade-off between scanning efficiency and imaging quality. Finally,
the promising potential of the tactile tomography system is
illustrated through its application in defect detection of
encapsulated FPCBs.

\section*{2.\ System design and instrumentation}

The configuration of the tactile tomography system is shown in
Fig.~1, and the photograph of the developed tactile tomography
system is presented in Fig.~2. The tactile tomography system
consists of a scanning module, a probe module, and a closed-loop
feedback control system.

\subsection*{2.1.\ Scanning module}

The mechanical design of the scanning module and probe module is
shown in Fig.~3. This design incorporates two linear actuators
(MBDLN25SE, Panasonic) and two servo motors (MSMF022L1U2M, Panasonic)
that together form an XY stage, allowing for precise adjustments of
the measurement position on the sample. A rodless cylinder
(CDHD-4D52AAP1) is utilized to drive the probe module, applying
pressure to the sample in the vertical direction, and a grating
sensor (RXF1188508-01) measures the depth of the applied pressure.
The movement speeds of the XY stage and Z stage can be adjusted
from 0.1 to 500 mm/min and from 0.1 to 400 mm/min, respectively.
The XY stage and Z stage achieve a limit resolution of step
displacement of 10 $\mu$m and

\end{document}
